\section{La revolución del post-entrenamiento: RLVR convierte la «recompensa verificable» en la vía principal}

\subsection{RLHF, RLVR y el post-training recipe completo}

Una de las partes más claras de la entrevista es la receta de post-entrenamiento: el mid-training aporta primero la estructura de la tarea, RLVR aplica aprendizaje por refuerzo con objetivos verificables, y RLHF hace al final la convergencia de estilo y usabilidad. Este orden explica por qué muchos modelos nuevos avanzan rápido en matemáticas y código, mientras mantienen controlada su «forma de hablar».

\begin{importantbox}{Por qué RLVR estalló en 2025--2026}
\begin{itemize}[leftmargin=1.5em]
    \item Tareas como matemáticas y código pueden verificarse automáticamente como correctas o incorrectas, con una señal de recompensa clara;
    \item la función de recompensa no depende de un reward model que «simula preferencias», lo que reduce el riesgo de reward hacking;
    \item en lo técnico permite paralelizar el rollout, lo que facilita seguir escalando el rendimiento del entrenamiento.
\end{itemize}
\end{importantbox}

\begin{figure}[H]
\centering
\includegraphics[width=0.88\textwidth]{frame-02.jpg}
\caption{Discusión concentrada en la entrevista sobre RLVR y la receta de post-entrenamiento\protect\footnotemark}
\end{figure}
\footnotetext{Intervalo de tiempo en el video: 01:47:10--01:47:20.}

\subsection{La arquitectura actor--learner y el costo de infraestructura}

Nathan explica el entrenamiento con RL como la colaboración de dos tipos de nodos de cómputo: el actor se encarga de generar y muestrear, y el learner de actualizar los gradientes. El primero puede estar geográficamente disperso, el segundo necesita una interconexión estrecha. Esta estructura difiere del modo de un solo clúster grande propio del pre-training tradicional, por lo que la planificación, la comunicación, la tolerancia a fallos y el retorno de datos resultan más complejos.

\begin{knowledgebox}{Por qué la ingeniería de RL suele «parecer más desordenada que la del preentrenamiento»}
\begin{itemize}[leftmargin=1.5em]
    \item Los datos se generan en línea, no se recorren directamente a partir de un corpus estático;
    \item la señal de retorno puede aparecer con retraso, lo que hace más lenta la retroalimentación del ajuste de parámetros;
    \item analizar los casos fallidos exige revisar tanto las trayectorias de la política como los registros del sistema;
    \item la inferencia y el entrenamiento están acoplados, y cualquier fluctuación en uno de los dos lados puede lastrar la eficiencia global.
\end{itemize}
\end{knowledgebox}

\begin{warningbox}{RLVR no es una máquina que «genere conocimiento nuevo automáticamente»}
La entrevista también mantiene con claridad una reserva: el RL se parece más a extraer y reordenar capacidades que el modelo ya posee, que a inyectar conocimiento nuevo de la nada. Si el corpus de preentrenamiento cubre de forma insuficiente cierto tipo de problemas, RLVR puede mejorar la calidad del proceso, pero difícilmente puede producir de la nada hechos que faltan. Este límite es clave para la planeación de la ruta a seguir.
\end{warningbox}

\subsection{Resumen de la sección}

La etapa de post-entrenamiento ha pasado de ser un «apéndice del ajuste fino» a convertirse en el campo de batalla principal. El auge de RLVR proviene de la combinación entre la verificabilidad de las tareas y la escalabilidad técnica, pero no sustituye a los datos de preentrenamiento de alta calidad.
